\documentclass[11pt]{article}

% Packages
\usepackage[a4paper,margin=1in]{geometry}
\usepackage{amsmath,amssymb}
\usepackage{enumitem}
\usepackage{fancyhdr}

% Header/Footer
\pagestyle{fancy}
\fancyhf{}
\lhead{CSE400: Fundamentals of Probability in Computing}
\rhead{Lecture 16 Scribe}
\cfoot{\thepage}

\begin{document}

\title{\textbf{Lecture 16: Joint Random Variables \& Joint Distributions}}
\author{}
\date{}
\maketitle

\section{Introduction: Why Study Joint Random Variables?}

\subsection{Motivation}
In many real-world systems, we are interested in analyzing two or more random variables simultaneously.

\subsection{Key Idea}
Let:
\[
X = \text{Random Variable 1}, \quad Y = \text{Random Variable 2}
\]
These variables may be correlated.

\subsection{Examples}

\textbf{Smart Transportation}
\begin{itemize}
    \item $X =$ Vehicle Speed
    \item $Y =$ Traffic Density
\end{itemize}

\textbf{Wireless Network Performance}
\begin{itemize}
    \item $X =$ Signal-to-Noise Ratio (SNR)
    \item $Y =$ Data Throughput
\end{itemize}

\textbf{Other Examples}
\begin{itemize}
    \item Weather Prediction: Rainfall, Temperature
    \item Drone Delivery: Wind Speed, Delivery Time
    \item Rolling Two Dice
\end{itemize}

\section{Discrete Case -- Joint PMF}

\subsection{Definition}
For discrete random variables $X$ and $Y$, the joint probability mass function (PMF) is:
\[
P(X = x, Y = y)
\]

\subsection{Interpretation}
It represents the probability of simultaneous outcomes of $X$ and $Y$.

\subsection{Joint PMF Table}
Values of $X$ and $Y$ are arranged in a table where each cell represents:
\[
P(X = x_i, Y = y_j)
\]

\section{Example: Rolling Two Dice}

\subsection{Setup}
Two dice are rolled.

\[
X = \text{Outcome of first die}, \quad Y = \text{Outcome of second die}
\]

\subsection{Objective}
Construct the joint PMF and evaluate joint events such as:
\[
X + Y = 7, \quad X > Y
\]

\section{Continuous Case -- Joint PDF}

\subsection{Definition}
For continuous random variables $X$ and $Y$, the joint probability density function (PDF) is:
\[
f_{X,Y}(x,y)
\]

\subsection{Interpretation}
Probability over a region $R$:
\[
P((X,Y) \in R) = \iint_R f_{X,Y}(x,y)\,dx\,dy
\]

\subsection{Geometric Interpretation}
Probability corresponds to volume under the surface over a region.

\section{Worked Example (Continuous Case)}

Tasks include:
\begin{itemize}
    \item Validating the PDF
    \item Computing probabilities over regions
    \item Performing double integration
\end{itemize}

\section{Joint Cumulative Distribution Function (Joint CDF)}

\subsection{Definition}
\[
F_{X,Y}(x,y) = P(X \leq x, Y \leq y)
\]

\section{Joint CDF -- Continuous Case}

Relation with joint PDF:
\[
F_{X,Y}(x,y) = \int_{-\infty}^{x} \int_{-\infty}^{y} f_{X,Y}(u,v)\,du\,dv
\]

\section{Joint CDF -- Discrete Case}

Computed using summation:
\[
F_{X,Y}(x,y) = \sum_{u \leq x} \sum_{v \leq y} P(X = u, Y = v)
\]

\section{Properties of the Joint CDF}

\begin{enumerate}
    \item \textbf{Non-decreasing Property:} Increasing in both $x$ and $y$
    \item \textbf{Limits:}
    \[
    F_{X,Y}(-\infty, y) = 0
    \]
    \[
    F_{X,Y}(x, -\infty) = 0
    \]
    \[
    F_{X,Y}(\infty, \infty) = 1
    \]
    \item \textbf{Right-continuity}
    \item \textbf{Probability Extraction:} Probabilities over regions can be derived using CDF values
\end{enumerate}

\section{Summary}

\begin{itemize}
    \item Joint random variables model multiple dependent quantities
    \item Two main cases:
    \begin{itemize}
        \item Discrete $\rightarrow$ Joint PMF
        \item Continuous $\rightarrow$ Joint PDF
    \end{itemize}
    \item Joint CDF provides cumulative probabilities
    \item Applications include:
    \begin{itemize}
        \item Transportation
        \item Wireless networks
        \item Weather systems
        \item Logistics
    \end{itemize}
\end{itemize}

\end{document}